\section{Rice theorems on the infinite line}\label{sec:line-results}
For the line categories, use the definitions of Section~\ref{sec:framework}
with spatial lattice $\Z$ and averaging intervals $[0,l)\cap\Z$.
In particular $E$ requires a unique stationary state among \emph{all}
states, attracting every initial state locally; extreme points are taken
in the jointly TI and stationary set. Use the same five-coordinate
notation $\sigma$ and the same maximal value $\sigma_\infty$.
The wavevector is in $\R/2\pi\Z$ and is nonzero for a Bragg label.

A classical line input has a nonempty finite alphabet, an arbitrary
finite stencil in $\Z$, a deterministic finite table, and independent
uniform whole-symbol replacement at rate $e$ after each update.
A quantum line input has finitely many qubits per site and a finite
repeating list of finite-range local CPTP gates with exact algebraic
entries and fresh product environments. Gates within a layer act on
disjoint internal tensor factors. After \emph{every} layer, independently
on every qubit, apply
\[
 \cD^{(1)}_e(X)=(1-e)X+e\Tr(X)I_2/2.
\]
This includes preparation, reset and idle layers. The full noisy round
is promised TI. The finite presentation, gate arity, number of wires,
range and period may vary with the input. Arbitrary entangled initial
states are allowed. Gate arity in the constructed family grows with
the program; no bounded-arity primitive-library assertion is made.
The general-site, once-per-round whole-site quantum category of
Theorem~\ref{thm:refined} is not treated on the line.

Fix the native interpreter, track order and tape template of
\cite[Section~3 and Appendix~A]{DMline}. Choose $w$ to be the least
admissible integer for the finite requirements of that construction,
including $w\ge10^{10}$. Put $P=5\cdot2^w$ and let $T$ be its fixed
native clock period. Both exceed $2w$. These choices are independent
of the program. Define
\begin{align}
 F_*&=\max(P,T),& M_*&=1000w(P+T+1)^3,\notag\\
 D_*&=8(P+2)^2(T+2)^2,&
 A_*&=(4M_*(P+1)+8)(4M_*(T+1)+8),\notag\\
 R_*&=4M_*F_*,& V_*&=10^9R_*^4A_*F_*^4,\notag\\
 e&=\frac{1}{4D_*w^2V_*^2}.&&\label{eq:line-rate}
\end{align}
These are definite computable integers and a definite rational once the
fixed interpreter is chosen. Only the existence of one such rate is
asserted; no decimal value is claimed. In particular $0<e<1/32$.

\begin{theorem}[Classical Rice theorem on the line]\label{thm:line-classical}
At the fixed rational rate \eqref{eq:line-rate}, every predicate of
$(N,E,L)$ nonconstant on the realized classical line data is undecidable.
\end{theorem}

\begin{theorem}[Quantum Rice theorem on the line]\label{thm:line-quantum}
For the line quantum inputs with independent single-qubit depolarization
after every layer specified above, at the same rate
\eqref{eq:line-rate}, every predicate of $(N,E,L)$ nonconstant on the
realized data is undecidable.
\end{theorem}

\begin{theorem}[Refined line theorem]\label{thm:line-refined}
In each line category separately, every predicate of
$\sigma=(N,N_{\ord},N_{\Br},E,L)$ nonconstant on the realized values is
undecidable at \eqref{eq:line-rate}. The value $\sigma_\infty$ is realized.
If a predicate is false there, halting many-one reduces to it; if true
there, nonhalting many-one reduces to it. These conclusions hold also
when every infinite count is reported as a single infinity symbol.
\end{theorem}

The proofs occupy Sections~\ref{sec:line-hierarchy}--\ref{sec:line-preservation}.
Compactness, extreme lifts and the implication from a nonzero Bragg
atom to covariance order follow by the full proofs of
Lemmas~\ref{lem:compact} and~\ref{lem:two-order}, with intervals in place
of cubes. A line circuit has a finite backward cone and independent
environments in distant cones, exactly as those proofs require.
For quantum changes of coordinates, the depolarizing convention is
covariant under products of single-qubit unitaries and wire permutations.
An entangling on-site conjugation need not remain in this input category.
